% Rana Neslihan Kurt, 2300006496
% Options for packages loaded elsewhere
\PassOptionsToPackage{unicode}{hyperref}
\PassOptionsToPackage{hyphens}{url}
\PassOptionsToPackage{dvipsnames,svgnames,x11names}{xcolor}
%
\documentclass[
  10pt,
]{article}
\usepackage{amsmath,amssymb}
\usepackage{iftex}
\ifPDFTeX
  \usepackage[T1]{fontenc}
  \usepackage[utf8]{inputenc}
  \usepackage{textcomp} % provide euro and other symbols
\else % if luatex or xetex
  \usepackage{unicode-math} % this also loads fontspec
  \defaultfontfeatures{Scale=MatchLowercase}
  \defaultfontfeatures[\rmfamily]{Ligatures=TeX,Scale=1}
\fi
\usepackage{lmodern}
\ifPDFTeX\else
  % xetex/luatex font selection
\fi
% Use upquote if available, for straight quotes in verbatim environments
\IfFileExists{upquote.sty}{\usepackage{upquote}}{}
\IfFileExists{microtype.sty}{% use microtype if available
  \usepackage[]{microtype}
  \UseMicrotypeSet[protrusion]{basicmath} % disable protrusion for tt fonts
}{}
\makeatletter
\@ifundefined{KOMAClassName}{% if non-KOMA class
  \IfFileExists{parskip.sty}{%
    \usepackage{parskip}
  }{% else
    \setlength{\parindent}{0pt}
    \setlength{\parskip}{6pt plus 2pt minus 1pt}}
}{% if KOMA class
  \KOMAoptions{parskip=half}}
\makeatother
\usepackage{xcolor}
\usepackage[margin=2.2cm]{geometry}
\usepackage{listings}
\newcommand{\passthrough}[1]{#1}
\lstset{defaultdialect=[5.3]Lua}
\lstset{defaultdialect=[x86masm]Assembler}
\usepackage{longtable,booktabs,array}
\usepackage{calc} % for calculating minipage widths
% Correct order of tables after \paragraph or \subparagraph
\usepackage{etoolbox}
\makeatletter
\patchcmd\longtable{\par}{\if@noskipsec\mbox{}\fi\par}{}{}
\makeatother
% Allow footnotes in longtable head/foot
\IfFileExists{footnotehyper.sty}{\usepackage{footnotehyper}}{\usepackage{footnote}}
\makesavenoteenv{longtable}
\setlength{\emergencystretch}{3em} % prevent overfull lines
\providecommand{\tightlist}{%
  \setlength{\itemsep}{0pt}\setlength{\parskip}{0pt}}
\setcounter{secnumdepth}{-\maxdimen} % remove section numbering
\usepackage{listings}
\usepackage{xcolor}
\lstset{basicstyle=\ttfamily\footnotesize,breaklines=true,breakatwhitespace=false,columns=fullflexible,keepspaces=true,frame=single,framerule=0.3pt,rulecolor=\color{gray!50},backgroundcolor=\color{gray!6},showstringspaces=false,literate={→}{{$\rightarrow$}}1 {—}{{---}}1 {…}{{\ldots}}1}
\usepackage{etoolbox}
\AtBeginEnvironment{longtable}{\footnotesize}
\setlength{\emergencystretch}{3em}
\usepackage{fvextra}
\DefineVerbatimEnvironment{Highlighting}{Verbatim}{breaklines,commandchars=\\\{\}}
\sloppy
\ifLuaTeX
  \usepackage{selnolig}  % disable illegal ligatures
\fi
\IfFileExists{bookmark.sty}{\usepackage{bookmark}}{\usepackage{hyperref}}
\IfFileExists{xurl.sty}{\usepackage{xurl}}{} % add URL line breaks if available
\urlstyle{same}
\hypersetup{
  colorlinks=true,
  linkcolor={Maroon},
  filecolor={Maroon},
  citecolor={Blue},
  urlcolor={Blue},
  pdfcreator={LaTeX via pandoc}}

\title{Destructuring Field Guide}
\author{Rana Neslihan Kurt, 2300006496}
\date{SEN0415 Advanced Web Programming, Assignment 4}

\begin{document}
\maketitle

{
\hypersetup{linkcolor=}
\setcounter{tocdepth}{3}
\tableofcontents
}
Name: Rana Neslihan Kurt\\
Student number: 2300006496

Bonuses claimed: \textbf{Bonus 1: LaTeX}, \textbf{Bonus 2: Depth},
\textbf{Bonus 3: Real code}, \textbf{Bonus 4: Git repository} (sections
at the end).

My data: my own TV show list --- \textbf{Suits}, \textbf{Awake} and
\textbf{The Mentalist}. Every output below is copied from
\passthrough{\lstinline!output.txt!} (a real run of
\passthrough{\lstinline!npm run all!} on my machine, Node 24.11.1,
Windows).

\hypertarget{run}{%
\subsection{Run}\label{run}}

\begin{lstlisting}
npm install
npm run all      # runs the JS and TS examples
npm run check    # type-checks in strict mode
\end{lstlisting}

\hypertarget{field-guide}{%
\subsection{Field guide}\label{field-guide}}

The idea behind all twelve forms: the left side is a \textbf{pattern}
that describes the shape of the value, and every name in the pattern
becomes a variable. Arrays match \textbf{by position}, objects match
\textbf{by name}. A missing item never throws, it becomes
\passthrough{\lstinline!undefined!}. Only destructuring
\passthrough{\lstinline!undefined!} or \passthrough{\lstinline!null!}
itself throws.

\hypertarget{form-1-values-by-position-const-x-y-f}{%
\subsubsection{\texorpdfstring{Form 1: values by position,
\texttt{const\ {[}x,\ y{]}\ =\ f()}}{Form 1: values by position, const {[}x, y{]} = f()}}\label{form-1-values-by-position-const-x-y-f}}

\textbf{Syntax:}
\passthrough{\lstinline!const [shortestShow, longestShow] = shortestAndLongest(shows)!}

\textbf{Example} (\passthrough{\lstinline!js/arrays.mjs!}):

\begin{lstlisting}
function shortestAndLongest(shows) {
  const sorted = [...shows].sort((a, b) => a.seasons - b.seasons)
  return [sorted[0].name, sorted[sorted.length - 1].name]
}
const [shortestShow, longestShow] = shortestAndLongest([
  { name: 'Suits', seasons: 9 }, { name: 'Awake', seasons: 1 }, { name: 'The Mentalist', seasons: 7 },
])
console.log('form 1 (my data):', shortestShow, longestShow)
\end{lstlisting}

\textbf{Output:} \passthrough{\lstinline!form 1 (my data): Awake Suits!}

\textbf{Without destructuring:}
\passthrough{\lstinline!const result = shortestAndLongest(shows); const shortestShow = result[0]; const longestShow = result[1]!}

\textbf{Use it when:} a function returns several values in a fixed order
(like \passthrough{\lstinline!useState!} returns
\passthrough{\lstinline![value, setter]!}).

\textbf{Watch out:} the order is everything. If I swap the names on the
left, I swap the meaning, and nothing warns me.

\hypertarget{form-2-skip-and-collect-const-second-...others-list}{%
\subsubsection{\texorpdfstring{Form 2: skip and collect,
\texttt{const\ {[},\ second,\ ...others{]}\ =\ list}}{Form 2: skip and collect, const {[}, second, ...others{]} = list}}\label{form-2-skip-and-collect-const-second-...others-list}}

\textbf{Syntax:}
\passthrough{\lstinline!const [, secondShow, ...otherShows] = myShows!}

\textbf{Example} (\passthrough{\lstinline!js/arrays.mjs!}):

\begin{lstlisting}
const myShows = ['Suits', 'Awake', 'The Mentalist']
const [, secondShow, ...otherShows] = myShows
console.log('form 2:', secondShow, otherShows)
const [, , , fourthShow, ...noMoreShows] = myShows
console.log('form 2 (too short):', fourthShow, noMoreShows)
\end{lstlisting}

\textbf{Output:}

\begin{lstlisting}
form 2: Awake [ 'The Mentalist' ]
form 2 (too short): undefined []
\end{lstlisting}

\textbf{Without destructuring:}
\passthrough{\lstinline!const secondShow = myShows[1]; const otherShows = myShows.slice(2)!}

\textbf{Use it when:} I need one item at a known place and the rest as a
list.

\textbf{Watch out:} the empty comma skips an item. Rest is
\textbf{always an array}: with one item left it is
\passthrough{\lstinline![ 'The Mentalist' ]!}, not a string, and with
nothing left it is \passthrough{\lstinline![]!}, not
\passthrough{\lstinline!undefined!}. Rest must be the last item in the
pattern.

\hypertarget{form-3-swap-a-b-b-a}{%
\subsubsection{\texorpdfstring{Form 3: swap,
\texttt{{[}a,\ b{]}\ =\ {[}b,\ a{]}}}{Form 3: swap, {[}a, b{]} = {[}b, a{]}}}\label{form-3-swap-a-b-b-a}}

\textbf{Syntax:}
\passthrough{\lstinline!;[favorite, secondFavorite] = [secondFavorite, favorite]!}

\textbf{Example} (\passthrough{\lstinline!js/arrays.mjs!}):

\begin{lstlisting}
let favorite = 'Suits'
let secondFavorite = 'Awake'
;[favorite, secondFavorite] = [secondFavorite, favorite]
console.log('form 3:', favorite, secondFavorite)
\end{lstlisting}

\textbf{Output:} \passthrough{\lstinline!form 3: Awake Suits!}

\textbf{Without destructuring:}
\passthrough{\lstinline!const temp = favorite; favorite = secondFavorite; secondFavorite = temp!}

\textbf{Use it when:} two variables must change places, without a
temporary variable.

\textbf{Watch out:} the variables must be \passthrough{\lstinline!let!},
and the line needs the \passthrough{\lstinline!;!} in front. Without it
JavaScript joins the line with the line above (see My own log, entry 2).

\hypertarget{form-4-array-defaults-const-a-1-b-2}{%
\subsubsection{\texorpdfstring{Form 4: array defaults,
\texttt{const\ {[}a\ =\ 1,\ b\ =\ 2{]}\ =\ {[}{]}}}{Form 4: array defaults, const {[}a = 1, b = 2{]} = {[}{]}}}\label{form-4-array-defaults-const-a-1-b-2}}

\textbf{Syntax:}
\passthrough{\lstinline!const [firstPick = 'Suits', secondPick = 'Awake'] = list!}

\textbf{Example} (\passthrough{\lstinline!js/arrays.mjs!}):

\begin{lstlisting}
const [firstPick = 'Suits', secondPick = 'Awake'] = []
const [pickA = 'Suits', pickB = 'Awake'] = ['The Mentalist', undefined]
const [pickC = 'Suits', pickD = 'Awake'] = [null, 'The Mentalist']
\end{lstlisting}

\textbf{Output:}

\begin{lstlisting}
form 4 (empty): Suits Awake
form 4 (undefined): The Mentalist Awake
form 4 (null): null The Mentalist
\end{lstlisting}

\textbf{Without destructuring:}
\passthrough{\lstinline!const pickC = list[0] === undefined ? 'Suits' : list[0]!}

\textbf{Use it when:} an item may be missing and there is a sensible
fallback.

\textbf{Watch out:} the default is used \textbf{only for
\passthrough{\lstinline!undefined!}}. \passthrough{\lstinline!null!} is
a real value, so \passthrough{\lstinline!pickC!} stays
\passthrough{\lstinline!null!}.

\hypertarget{form-5-properties-by-name-let-a-b-t}{%
\subsubsection{\texorpdfstring{Form 5: properties by name,
\texttt{let\ \{\ a,\ b\ \}\ =\ t}}{Form 5: properties by name, let \{ a, b \} = t}}\label{form-5-properties-by-name-let-a-b-t}}

\textbf{Syntax:}
\passthrough{\lstinline!const \{ showTitle, mainCharacter \} = favoriteShow!}

\textbf{Example} (\passthrough{\lstinline!js/objects.mjs!}):

\begin{lstlisting}
const favoriteShow = { showTitle: 'The Mentalist', mainCharacter: 'Patrick Jane', seasons: 7 }
const { showTitle, mainCharacter } = favoriteShow
console.log('form 5 (my data):', showTitle, mainCharacter)
\end{lstlisting}

\textbf{Output:}
\passthrough{\lstinline!form 5 (my data): The Mentalist Patrick Jane!}

\textbf{Without destructuring:}
\passthrough{\lstinline!const showTitle = favoriteShow.showTitle; const mainCharacter = favoriteShow.mainCharacter!}

\textbf{Use it when:} I need a few fields of an object and want short
names for them.

\textbf{Watch out:} the variable name must be exactly the property name.
A typo gives \passthrough{\lstinline!undefined!}, not an error.

\hypertarget{form-6-rename-const-a-alpha-t}{%
\subsubsection{\texorpdfstring{Form 6: rename,
\texttt{const\ \{\ a:\ alpha\ \}\ =\ t}}{Form 6: rename, const \{ a: alpha \} = t}}\label{form-6-rename-const-a-alpha-t}}

\textbf{Syntax:}
\passthrough{\lstinline!const \{ name: showName, seasons: seasonCount \} = show!}

\textbf{Example} (\passthrough{\lstinline!js/objects.mjs!}):

\begin{lstlisting}
const show = { name: 'Suits', network: 'USA Network', seasons: 9 }
const { name: showName, seasons: seasonCount } = show
console.log('form 6:', showName, seasonCount)
try { console.log(name) } catch (error) { console.log('form 6 (old name):', error.message) }
\end{lstlisting}

\textbf{Output:}

\begin{lstlisting}
form 6: Suits 9
form 6 (old name): name is not defined
\end{lstlisting}

\textbf{Without destructuring:}
\passthrough{\lstinline!const showName = show.name; const seasonCount = show.seasons!}

\textbf{Use it when:} the property name clashes with a variable I
already have, or is not clear here.

\textbf{Watch out:} left of the colon is the \textbf{source}, right is
the \textbf{new variable}. The old name \passthrough{\lstinline!name!}
is not a variable afterwards.

\hypertarget{form-7-object-defaults-const-a-1-t}{%
\subsubsection{\texorpdfstring{Form 7: object defaults,
\texttt{const\ \{\ a\ =\ 1\ \}\ =\ t}}{Form 7: object defaults, const \{ a = 1 \} = t}}\label{form-7-object-defaults-const-a-1-t}}

\textbf{Syntax:}
\passthrough{\lstinline!const \{ network = 'unknown network', rating = 'no rating' \} = awake!}

\textbf{Example} (\passthrough{\lstinline!js/objects.mjs!}):

\begin{lstlisting}
const awake = { name: 'Awake', seasons: 1, rating: undefined }
const { network = 'unknown network', rating = 'no rating' } = awake
const mentalist = { name: 'The Mentalist', network: null }
const { network: mentalistNetwork = 'unknown network' } = mentalist
\end{lstlisting}

\textbf{Output:}

\begin{lstlisting}
form 7: unknown network no rating
form 7 (null): null
\end{lstlisting}

\textbf{Without destructuring:}
\passthrough{\lstinline!const network = awake.network === undefined ? 'unknown network' : awake.network!}

\textbf{Use it when:} a field is optional (settings, options objects,
props).

\textbf{Watch out:} a missing field and a field set to
\passthrough{\lstinline!undefined!} both get the default,
\passthrough{\lstinline!null!} does not. Rename and default can be
combined:
\passthrough{\lstinline!\{ network: mentalistNetwork = '...' \}!}.

\hypertarget{form-8-rest-in-an-object-const-a-...others-t}{%
\subsubsection{\texorpdfstring{Form 8: rest in an object,
\texttt{const\ \{\ a,\ ...others\ \}\ =\ t}}{Form 8: rest in an object, const \{ a, ...others \} = t}}\label{form-8-rest-in-an-object-const-a-...others-t}}

\textbf{Syntax:}
\passthrough{\lstinline!const \{ id, ...showInfo \} = suits!}

\textbf{Example} (\passthrough{\lstinline!js/objects.mjs!}):

\begin{lstlisting}
const suits = { id: 's1', name: 'Suits', network: 'USA Network', seasons: 9 }
const { id, ...showInfo } = suits
console.log('form 8:', id, showInfo)
const updatedInfo = { ...showInfo, seasons: 10 }
console.log('form 8 (spread):', updatedInfo)
\end{lstlisting}

\textbf{Output:}

\begin{lstlisting}
form 8: s1 { name: 'Suits', network: 'USA Network', seasons: 9 }
form 8 (spread): { name: 'Suits', network: 'USA Network', seasons: 10 }
\end{lstlisting}

\textbf{Without destructuring:}
\passthrough{\lstinline!const id = suits.id; const showInfo = Object.assign(\{\}, suits); delete showInfo.id!}

\textbf{Use it when:} I want to remove one field and keep the rest (very
common with React props, see P4).

\textbf{Watch out:} the same three dots do opposite jobs. On the
\textbf{left} of \passthrough{\lstinline!=!} it is \textbf{rest}
(collect into a new object). On the \textbf{right} it is \textbf{spread}
(expand an object into a new one). With spread, a field written later
overwrites an earlier one (\passthrough{\lstinline!seasons!} 9 → 10).

\hypertarget{form-9-nested-patterns-const-p-x-list-first-t}{%
\subsubsection{\texorpdfstring{Form 9: nested patterns,
\texttt{const\ \{\ p:\ \{\ x\ \},\ list:\ {[}first{]}\ \}\ =\ t}}{Form 9: nested patterns, const \{ p: \{ x \}, list: {[}first{]} \} = t}}\label{form-9-nested-patterns-const-p-x-list-first-t}}

\textbf{Syntax:}
\passthrough{\lstinline!const \{ owner: \{ firstName \}, shows: [firstShow, ...laterShows] \} = myWatchlist!}

\textbf{Example} (\passthrough{\lstinline!js/objects.mjs!}):

\begin{lstlisting}
const myWatchlist = {
  owner: { firstName: 'Rana', favoriteGenre: 'drama' },
  shows: ['Suits', 'Awake', 'The Mentalist'],
}
const { owner: { firstName }, shows: [firstShow, ...laterShows] } = myWatchlist
try {
  const { creator: { firstName: creatorName } } = myWatchlist
} catch (error) { console.log('form 9 (missing nested):', error.message) }
\end{lstlisting}

\textbf{Output:}

\begin{lstlisting}
form 9: Rana Suits [ 'Awake', 'The Mentalist' ]
form 9 (missing nested): Cannot read properties of undefined (reading 'firstName')
\end{lstlisting}

\textbf{Without destructuring:}
\passthrough{\lstinline!const firstName = myWatchlist.owner.firstName; const firstShow = myWatchlist.shows[0]; const laterShows = myWatchlist.shows.slice(1)!}

\textbf{Use it when:} I need values from deep inside a structure, for
example an API response.

\textbf{Watch out:} \passthrough{\lstinline!owner!} and
\passthrough{\lstinline!shows!} are only the path, \textbf{not
variables}. If the outer object is missing, it throws; a default like
\passthrough{\lstinline!\{ creator: \{ firstName \} = \{\} \}!} prevents
it.

\hypertarget{form-10-computed-key-const-key-value-t}{%
\subsubsection{\texorpdfstring{Form 10: computed key,
\texttt{const\ \{\ {[}key{]}:\ value\ \}\ =\ t}}{Form 10: computed key, const \{ {[}key{]}: value \} = t}}\label{form-10-computed-key-const-key-value-t}}

\textbf{Syntax:}
\passthrough{\lstinline!const \{ [pickedShow]: pickedRating \} = myRatings!}

\textbf{Example} (\passthrough{\lstinline!js/objects.mjs!}):

\begin{lstlisting}
const myRatings = { Suits: 9, Awake: 8, 'The Mentalist': 8.5 }
const pickedShow = 'Awake'
const { [pickedShow]: pickedRating } = myRatings
const missingKey = 'Lost'
const { [missingKey]: missingRating = 'not rated' } = myRatings
\end{lstlisting}

\textbf{Output:}

\begin{lstlisting}
form 10: Awake 8
form 10 (missing key): not rated
\end{lstlisting}

\textbf{Without destructuring:}
\passthrough{\lstinline!const pickedRating = myRatings[pickedShow]!}

\textbf{Use it when:} the name of the field is only known at run time
(stored in a variable).

\textbf{Watch out:} a new name after the colon is required, because
JavaScript cannot make a variable name out of a value. The square
brackets mean ``use the value of \passthrough{\lstinline!pickedShow!}'',
not the word \passthrough{\lstinline!pickedShow!}.

\hypertarget{form-11-patterns-in-parameters-function-f-a-1-b-options}{%
\subsubsection{\texorpdfstring{Form 11: patterns in parameters,
\texttt{function\ f(\{\ a\ =\ 1,\ b\ \}:\ Options\ =\ \{\})}}{Form 11: patterns in parameters, function f(\{ a = 1, b \}: Options = \{\})}}\label{form-11-patterns-in-parameters-function-f-a-1-b-options}}

\textbf{Syntax:}
\passthrough{\lstinline!function describeShow(\{ title, seasons = 1, finished = true \} = \{\})!}

\textbf{Example} (\passthrough{\lstinline!js/params.mjs!}):

\begin{lstlisting}
function describeShow({ title, seasons = 1, finished = true } = {}) {
  return `${title}: ${seasons} season(s), finished: ${finished}`
}
console.log('form 11 (my data):', describeShow({ title: 'Awake' }))
console.log('form 11 (no argument):', describeShow())
function describeShowNoDefault({ title }) { return title }
try { describeShowNoDefault() } catch (error) { console.log('form 11 (without = {}):', error.message) }
const formatEpisode = ([season, episode]) => `S${season}E${episode}`
console.log('form 11b:', formatEpisode([2, 5]))
\end{lstlisting}

\textbf{Output:}

\begin{lstlisting}
form 11 (my data): Awake: 1 season(s), finished: true
form 11 (no argument): undefined: 1 season(s), finished: true
form 11 (without = {}): Cannot destructure property 'title' of 'undefined' as it is undefined.
form 11b: S2E5
\end{lstlisting}

\textbf{Without destructuring:}
\passthrough{\lstinline!function describeShow(options) \{ options = options || \{\}; const title = options.title; ... \}!}

\textbf{Use it when:} a function takes an options object. The call is
readable (\passthrough{\lstinline!describeShow(\{ title: 'Awake' \})!})
and the order of options does not matter. Every React component does
this with props.

\textbf{Watch out:} there are two defaults. The
\passthrough{\lstinline!= 1!} inside is for a missing \textbf{field}.
The \passthrough{\lstinline!= \{\}!} at the end is for a missing
\textbf{argument}. Without \passthrough{\lstinline!= \{\}!} the call
with no argument throws.

\hypertarget{form-12-patterns-in-loops-for-const-k-v-of-object.entrieso}{%
\subsubsection{\texorpdfstring{Form 12: patterns in loops,
\texttt{for\ (const\ {[}k,\ v{]}\ of\ Object.entries(o))}}{Form 12: patterns in loops, for (const {[}k, v{]} of Object.entries(o))}}\label{form-12-patterns-in-loops-for-const-k-v-of-object.entrieso}}

\textbf{Syntax:}
\passthrough{\lstinline!for (const [show, rating] of Object.entries(myRatings))!}

\textbf{Example} (\passthrough{\lstinline!js/params.mjs!}):

\begin{lstlisting}
const myRatings = { Suits: 9, Awake: 8, 'The Mentalist': 8.5 }
for (const [show, rating] of Object.entries(myRatings)) console.log('form 12 (object):', show, rating)
const lastWatched = new Map([['Suits', 'S9E10'], ['Awake', 'S1E13'], ['The Mentalist', 'S7E13']])
for (const [show, episode] of lastWatched) console.log('form 12 (map):', show, episode)
\end{lstlisting}

\textbf{Output:}

\begin{lstlisting}
form 12 (object): Suits 9
form 12 (object): Awake 8
form 12 (object): The Mentalist 8.5
form 12 (map): Suits S9E10
form 12 (map): Awake S1E13
form 12 (map): The Mentalist S7E13
\end{lstlisting}

\textbf{Without destructuring:}
\passthrough{\lstinline!for (const entry of Object.entries(myRatings)) \{ const show = entry[0]; const rating = entry[1] \}!}

\textbf{Use it when:} I loop over key--value pairs of an object or a
\passthrough{\lstinline!Map!}.

\textbf{Watch out:} \passthrough{\lstinline!Object.entries!} turns the
object into \passthrough{\lstinline![key, value]!} pairs first; a
\passthrough{\lstinline!Map!} already gives pairs. Every round of the
loop is a new Form 1 on one pair.

\hypertarget{pitfalls}{%
\subsection{Pitfalls}\label{pitfalls}}

All five are in \passthrough{\lstinline!js/pitfalls.mjs!}.

\textbf{Pitfall 1: destructuring \passthrough{\lstinline!undefined!} or
\passthrough{\lstinline!null!} throws.}
\passthrough{\lstinline!const \{ a \} = undefined!} →
\passthrough{\lstinline!pitfall 1: TypeError!}. \emph{Why:} to read a
property, JavaScript first needs an object.
\passthrough{\lstinline!undefined!} and \passthrough{\lstinline!null!}
have no properties at all, so there is nothing to read from. A missing
\emph{field} is fine (\passthrough{\lstinline!undefined!}), a missing
\emph{object} is not. Fix: a default like
\passthrough{\lstinline!= \{\}!} (Form 11).

\textbf{Pitfall 2: a default only works for
\passthrough{\lstinline!undefined!}, not for
\passthrough{\lstinline!null!}.}

\begin{lstlisting}
const suitsInfo = { network: null, seasons: undefined }
const { network = 'unknown', seasons = 1 } = suitsInfo
\end{lstlisting}

Output: \passthrough{\lstinline!pitfall 2: null 1!} \emph{Why:} the
default is checked with \passthrough{\lstinline!=== undefined!}.
\passthrough{\lstinline!null!} means ``a value that is empty on
purpose'', so JavaScript keeps it.

\textbf{Pitfall 3: assigning without
\passthrough{\lstinline!const!}/\passthrough{\lstinline!let!} needs
parentheses.}

\begin{lstlisting}
let currentShow
eval("{ currentShow } = { currentShow: 'Awake' }")          // SyntaxError
;({ currentShow, currentSeason = 1 } = { currentShow: 'Awake' }) // works
\end{lstlisting}

Output:

\begin{lstlisting}
pitfall 3 (no parentheses): SyntaxError - Unexpected token '='
pitfall 3 (with parentheses): Awake 1
\end{lstlisting}

\emph{Why:} a line that starts with \passthrough{\lstinline!\{!} is read
as a \textbf{block of code}, not as an object pattern, so the
\passthrough{\lstinline!=!} after it makes no sense. Parentheses force
JavaScript to read it as an expression. (I used
\passthrough{\lstinline!eval!} only so that the SyntaxError could be
caught and printed; otherwise the whole file would not start.)

\textbf{Pitfall 4: destructuring copies references, not objects.}

\begin{lstlisting}
const watchlist = { owner: 'Rana', shows: ['Suits', 'Awake'] }
const { shows } = watchlist
shows.push('The Mentalist')
const copy = { ...watchlist }
copy.shows.push('House')
copy.owner = 'Someone else'
\end{lstlisting}

Output:

\begin{lstlisting}
pitfall 4: [ 'Suits', 'Awake', 'The Mentalist' ]
pitfall 4 (spread copy): {
  owner: 'Rana',
  shows: [ 'Suits', 'Awake', 'The Mentalist', 'House' ]
}
\end{lstlisting}

\emph{Why:} \passthrough{\lstinline!shows!} is a second name for
\textbf{the same array}, so changing it changes
\passthrough{\lstinline!watchlist.shows!}. Even spread is a
\textbf{shallow} copy: \passthrough{\lstinline!owner!} was copied (so it
stayed \passthrough{\lstinline!'Rana'!}), but the inner array is still
shared.

\textbf{Pitfall 5: in objects the order does not matter, in arrays it
does.}

\begin{lstlisting}
const { seasons: suitsSeasons, title } = { title: 'Suits', seasons: 9 }
const [nameFirst, numberFirst] = ['Suits', 9]
const [numberNow, nameNow] = [9, 'Suits']
\end{lstlisting}

Output:

\begin{lstlisting}
pitfall 5 (object): Suits 9
pitfall 5 (array): Suits 9 | 9 Suits
\end{lstlisting}

\emph{Why:} object patterns match \textbf{by name}, array patterns match
\textbf{by position}. Writing the array pattern in the wrong order gives
wrong values without any error.

\hypertarget{typescript-topics-t1-to-t5}{%
\subsection{TypeScript topics (T1 to
T5)}\label{typescript-topics-t1-to-t5}}

All in \passthrough{\lstinline!ts/types.ts!}.

\textbf{T1: in a pattern, \passthrough{\lstinline!\{ name: string \}!}
is a rename, not a type.} I wrote
\passthrough{\lstinline!function showLabel(\{ name: string, seasons: number \})!}
and got:

\begin{lstlisting}
error TS7031: Binding element 'string' implicitly has an 'any' type.
error TS7031: Binding element 'number' implicitly has an 'any' type.
error TS2304: Cannot find name 'seasons'.
\end{lstlisting}

TypeScript made variables called \passthrough{\lstinline!string!} and
\passthrough{\lstinline!number!} (Form 6). The fix: the type goes
\textbf{after the whole pattern}:
\passthrough{\lstinline!function showLabel(\{ name, seasons \}: \{ name: string; seasons: number \}): string!}
→ \passthrough{\lstinline!T1 fixed: Suits (9)!}.

\textbf{T2: tuples and \passthrough{\lstinline!as const!}.}
\passthrough{\lstinline!function bestShow(): [string, number]!} gives
\passthrough{\lstinline!bestName: string!} and
\passthrough{\lstinline!bestRating: number!}, so
\passthrough{\lstinline!bestName.toUpperCase()!} and
\passthrough{\lstinline!bestRating + 1!} both compile
(\passthrough{\lstinline!T2 tuple: SUITS 10!}). A plain
\passthrough{\lstinline!['Awake', 1]!} is only
\passthrough{\lstinline!(string | number)[]!}, so each destructured item
is \passthrough{\lstinline!string | number | undefined!}, and
\passthrough{\lstinline!const seasonsAsNumber: number = looseSeasons!}
gives
\passthrough{\lstinline!error TS2322: Type 'string | number | undefined' is not assignable to type 'number'.!}
With \passthrough{\lstinline!as const!} the array becomes
\passthrough{\lstinline!readonly ['Awake', 1]!}:
\passthrough{\lstinline!fixedName!} has type
\passthrough{\lstinline!"Awake"!} and
\passthrough{\lstinline!fixedSeasons!} has type
\passthrough{\lstinline!1!}, so the same assignment works.

\textbf{T3: \passthrough{\lstinline!noUncheckedIndexedAccess!}.} With
this option \passthrough{\lstinline!const [firstToday] = watchedToday!}
(a \passthrough{\lstinline!string[]!}) has type
\passthrough{\lstinline!string | undefined!}, because the array could be
empty. \passthrough{\lstinline!firstToday.toUpperCase()!} gives
\passthrough{\lstinline!error TS18048: 'firstToday' is possibly 'undefined'.!}
Three fixes, all in the file: check with
\passthrough{\lstinline"if (firstToday !== undefined)"}, give a default
\passthrough{\lstinline![firstOrNone = 'nothing']!}, or use a tuple type
\passthrough{\lstinline![string, string]!} so TypeScript knows the item
exists. Output: \passthrough{\lstinline!T3 checked: SUITS!},
\passthrough{\lstinline!T3 default: SUITS!},
\passthrough{\lstinline!T3 tuple: SUITS!}.

\textbf{T4: optional fields and defaults.} In
\passthrough{\lstinline!(\{ title, rating, subtitles = 'English' \}: ShowSettings)!}
with \passthrough{\lstinline!rating?: number; subtitles?: string!}:
\passthrough{\lstinline!title!} is \passthrough{\lstinline!string!}
(required), \passthrough{\lstinline!rating!} is
\passthrough{\lstinline!number | undefined!} (optional, no default),
\passthrough{\lstinline!subtitles!} is \passthrough{\lstinline!string!}
(optional, but the default removes \passthrough{\lstinline!undefined!}).
Output: \passthrough{\lstinline!T4 all given: Suits | 9.0 | Turkish!}
and
\passthrough{\lstinline!T4 only title: Awake | no rating | English!}.

\textbf{T5: a compile error caused by destructuring.}
\passthrough{\lstinline!const \{ rating: mentalistRating \} = mentalistShow!}
where the object has no \passthrough{\lstinline!rating!}:
\passthrough{\lstinline!error TS2339: Property 'rating' does not exist on type '\{ title: string; seasons: number; \}'.!}
This is the only line with
\passthrough{\lstinline!// @ts-expect-error!}. At run time the value is
simply \passthrough{\lstinline!undefined!}
(\passthrough{\lstinline!T5 at runtime: undefined!}), which shows that
TypeScript catches the mistake that plain JavaScript lets through. Fix:
add \passthrough{\lstinline!rating?: number!} to the type and a default
→ \passthrough{\lstinline!T5 fixed: 0!}.

\hypertarget{props-and-children}{%
\subsection{Props and children}\label{props-and-children}}

Code: \passthrough{\lstinline!js/props.jsx!} (P1 to P6, run with tsx)
and \passthrough{\lstinline!ts/props.tsx!} (P1 to P7 with types).
Rendered with \passthrough{\lstinline!renderToStaticMarkup!}.

\hypertarget{p1-props-is-one-object}{%
\subsubsection{P1: props is one object}\label{p1-props-is-one-object}}

\begin{lstlisting}
function ShowCardWhole(props) {
  console.log('P1 props object:', props)
  return <article><h3>{props.name}</h3><p>{props.children}</p></article>
}
function ShowCardParts({ name, children }) {
  return <article><h3>{name}</h3><p>{children}</p></article>
}
\end{lstlisting}

Output:

\begin{lstlisting}
P1 props object: { name: 'Suits', children: 'A lawyer drama' }
P1 whole: <article><h3>Suits</h3><p>A lawyer drama</p></article>
P1 parts: <article><h3>Suits</h3><p>A lawyer drama</p></article>
\end{lstlisting}

A component gets \textbf{one} argument. The printed object shows that
the attribute \passthrough{\lstinline!name!} and the content between the
tags (\passthrough{\lstinline!children!}) are both fields of that one
object. \passthrough{\lstinline!(\{ name, children \})!} is \textbf{Form
5} (take fields by name) used in a \textbf{parameter} (\textbf{Form
11}); the result is the same as reading
\passthrough{\lstinline!props.name!} and
\passthrough{\lstinline!props.children!}.

\hypertarget{p2-children-is-only-a-prop}{%
\subsubsection{P2: children is only a
prop}\label{p2-children-is-only-a-prop}}

\begin{lstlisting}
<ShowCardParts name="Awake">Two worlds</ShowCardParts>
<ShowCardParts name="Awake" children="Two worlds" />
<ShowCardParts name="Awake" children="from attribute">from tags</ShowCardParts>
\end{lstlisting}

Output:

\begin{lstlisting}
P2 between tags: <article><h3>Awake</h3><p>Two worlds</p></article>
P2 as attribute: <article><h3>Awake</h3><p>Two worlds</p></article>
P2 both: <article><h3>Awake</h3><p>from tags</p></article>
\end{lstlisting}

The first two are identical: content between tags is just a shorter way
to write the \passthrough{\lstinline!children!} prop. When I gave both,
the content between the tags \textbf{won}, because JSX sets the
attributes first and then puts the content into
\passthrough{\lstinline!children!}, overwriting it.

\hypertarget{p3-what-children-can-be}{%
\subsubsection{P3: what children can be}\label{p3-what-children-can-be}}

Component:
\passthrough{\lstinline!function Box(\{ children \}) \{ return <div>\{children\}</div> \}!}

\begin{longtable}[]{@{}
  >{\raggedright\arraybackslash}p{(\columnwidth - 4\tabcolsep) * \real{0.3333}}
  >{\raggedright\arraybackslash}p{(\columnwidth - 4\tabcolsep) * \real{0.3333}}
  >{\raggedright\arraybackslash}p{(\columnwidth - 4\tabcolsep) * \real{0.3333}}@{}}
\toprule\noalign{}
\begin{minipage}[b]{\linewidth}\raggedright
children
\end{minipage} & \begin{minipage}[b]{\linewidth}\raggedright
rendered
\end{minipage} & \begin{minipage}[b]{\linewidth}\raggedright
appears?
\end{minipage} \\
\midrule\noalign{}
\endhead
\bottomrule\noalign{}
\endlastfoot
\passthrough{\lstinline!The Mentalist!} (string) &
\passthrough{\lstinline!<div>The Mentalist</div>!} & yes \\
\passthrough{\lstinline!\{7\}!} & \passthrough{\lstinline!<div>7</div>!}
& yes \\
\passthrough{\lstinline!\{0\}!} & \passthrough{\lstinline!<div>0</div>!}
& \textbf{yes} (0 is a number, not ``empty'') \\
\passthrough{\lstinline!<b>Suits</b>!} &
\passthrough{\lstinline!<div><b>Suits</b></div>!} & yes \\
\passthrough{\lstinline!<b>Suits</b><i>Awake</i>!} &
\passthrough{\lstinline!<div><b>Suits</b><i>Awake</i></div>!} & yes \\
\passthrough{\lstinline!\{['Suits', 'Awake']\}!} &
\passthrough{\lstinline!<div>SuitsAwake</div>!} & yes, glued together \\
\passthrough{\lstinline!\{undefined\}!},
\passthrough{\lstinline!\{null\}!}, \passthrough{\lstinline!\{false\}!},
\passthrough{\lstinline!\{true\}!}, no children &
\passthrough{\lstinline!<div></div>!} & no \\
\passthrough{\lstinline!\{() => 'Patrick Jane'\}!} &
\passthrough{\lstinline!<div></div>!} + warning \emph{Functions are not
valid as a React child} & no \\
function, and the component \textbf{calls} it:
\passthrough{\lstinline!children('Patrick Jane')!} &
\passthrough{\lstinline!<div><b>Patrick Jane</b></div>!} & yes \\
\passthrough{\lstinline!\{\{ name: 'Suits' \}\}!} (plain object) &
error: \emph{Objects are not valid as a React child (found: object with
keys \{name\})} & crashes \\
\end{longtable}

The surprise: \passthrough{\lstinline!false!} and
\passthrough{\lstinline!null!} disappear but \passthrough{\lstinline!0!}
is printed. That is why
\passthrough{\lstinline!\{count \&\& <List />\}!} can show a stray
\passthrough{\lstinline!0!} on the screen. A function is fine as
children only if the component treats it as a function and calls it.

\hypertarget{p4-separating-children-from-the-other-props}{%
\subsubsection{P4: separating children from the other
props}\label{p4-separating-children-from-the-other-props}}

\begin{lstlisting}
function Panel({ title, children, ...rest }) {
  console.log('P4 rest:', rest)
  return <section {...rest}><h2>{title}</h2>{children}</section>
}
<Panel title="My watchlist" id="list" className="dark" data-count="3">3 shows</Panel>
\end{lstlisting}

Output:

\begin{lstlisting}
P4 rest: { id: 'list', className: 'dark', 'data-count': '3' }
P4: <section id="list" class="dark" data-count="3"><h2>My watchlist</h2>3 shows</section>
\end{lstlisting}

\passthrough{\lstinline!title!} and \passthrough{\lstinline!children!}
are used by the component; \textbf{Form 8} collects everything else into
\passthrough{\lstinline!rest!}, and the spread
\passthrough{\lstinline!\{...rest\}!} passes it on to the real
\passthrough{\lstinline!<section>!}.

\textbf{Order decides who wins:}

\begin{lstlisting}
<section {...rest} className="fixed">   // P4 spread first: <section class="fixed">Suits</section>
<section className="fixed" {...rest}>   // P4 spread last:  <section class="mine">Suits</section>
\end{lstlisting}

(called with \passthrough{\lstinline!className="mine"!}). Like object
spread in Form 8, whatever is written \textbf{later} overwrites. Spread
first = the component has the last word; spread last = the caller can
override.

\textbf{\passthrough{\lstinline!key!} is not a prop.} With
\passthrough{\lstinline!<WithKey key="s1" name="Suits" />!},
\passthrough{\lstinline!props.key!} is
\passthrough{\lstinline!undefined!} and React warns:
\emph{\passthrough{\lstinline!key!} is not a prop. Trying to access it
will result in \passthrough{\lstinline!undefined!} being returned.}
React keeps \passthrough{\lstinline!key!} for itself (to track list
items).

\textbf{Props cannot be changed.}
\passthrough{\lstinline!props.show = 'changed'!} gave
\passthrough{\lstinline!TypeError - Cannot assign to read only property 'show' of object '\#<Object>'!}.
React freezes the props object; a component must treat props as
read-only.

\hypertarget{p5-props-or-children}{%
\subsubsection{P5: props or children?}\label{p5-props-or-children}}

\begin{lstlisting}
function ShowPageProps({ header, body, footer }) {
  return <main><header>{header}</header><div>{body}</div><footer>{footer}</footer></main>
}
function ShowPageChildren({ children }) {
  return <main>{children}</main>
}
\end{lstlisting}

Output:

\begin{lstlisting}
P5 props: <main><header><h1>Suits</h1></header><div>Harvey and Mike</div><footer><small>9 seasons</small></footer></main>
P5 children: <main><h1>Suits</h1><p>Harvey and Mike</p><small>9 seasons</small></main>
\end{lstlisting}

With \textbf{named props} the component decides the layout: the header
always goes into \passthrough{\lstinline!<header>!}, the footer into
\passthrough{\lstinline!<footer>!}, and the caller only fills the slots.
With \textbf{children} the component is just a frame and the caller
decides everything inside (here the
\passthrough{\lstinline!<header>!}/\passthrough{\lstinline!<footer>!}
tags are gone). I would choose \textbf{named props} when there are
several fixed places (a card with a title and actions, a page layout)
and \textbf{children} when there is one free area whose content the
component does not need to know (a panel, a modal, a button label).

\hypertarget{p6-defaults}{%
\subsubsection{P6: defaults}\label{p6-defaults}}

\begin{lstlisting}
function Rating({ show, stars = 'not rated yet' }) { return <p>{show}: {stars}</p> }
\end{lstlisting}

Output:

\begin{lstlisting}
P6 given: <p>Suits: 9</p>
P6 missing: <p>Awake: not rated yet</p>
P6 undefined: <p>Awake: not rated yet</p>
P6 null: <p>The Mentalist: </p>
\end{lstlisting}

This is \textbf{Form 7} in a parameter. A missing prop and
\passthrough{\lstinline!stars=\{undefined\}!} both get the default.
\passthrough{\lstinline!stars=\{null\}!} does \textbf{not}, and React
renders \passthrough{\lstinline!null!} as nothing, so the screen shows
\passthrough{\lstinline!The Mentalist:!} with an empty space. Same rule
as Form 4 and Pitfall 2.

\hypertarget{p7-types-tsprops.tsx}{%
\subsubsection{\texorpdfstring{P7: types
(\texttt{ts/props.tsx})}{P7: types (ts/props.tsx)}}\label{p7-types-tsprops.tsx}}

\begin{lstlisting}
type InfoProps = { label: string; children?: ReactNode }               // 1) by hand
type BadgeProps = PropsWithChildren<{ color: string }>                 // 2) adds children?: ReactNode
type WatchButtonProps = ComponentProps<'button'> & { show: string }    // 3) all props of a real <button>
function WatchButton({ show, children, ...buttonProps }: WatchButtonProps) {
  return <button {...buttonProps}>{children} {show}</button>
}
\end{lstlisting}

Output:

\begin{lstlisting}
P7 ReactNode: <p><b>Network:</b> USA Network</p>
P7 PropsWithChildren: <span class="gold">Favorite</span>
P7 ComponentProps: <button type="button" disabled="">Watch Suits</button>
\end{lstlisting}

\begin{itemize}
\tightlist
\item
  \passthrough{\lstinline!children?: ReactNode!} allows everything React
  can render (string, number, element, array,
  \passthrough{\lstinline!null!}, \passthrough{\lstinline!undefined!},
  boolean).
\item
  \passthrough{\lstinline!PropsWithChildren<T>!} is a shortcut that adds
  exactly that line to my type \passthrough{\lstinline!T!}.
\item
  \passthrough{\lstinline!ComponentProps<'button'>!} gives every
  attribute a real \passthrough{\lstinline!<button>!} accepts
  (\passthrough{\lstinline!type!}, \passthrough{\lstinline!disabled!},
  \passthrough{\lstinline!onClick!}, \ldots), so
  \passthrough{\lstinline!...buttonProps!} is typed and can be spread
  onto the button.
\item
  A function as children needs its own type:
  \passthrough{\lstinline!children: (name: string) => ReactNode!}.
\end{itemize}

What the compiler rejects (each line tried, real errors):

\begin{lstlisting}
<Info>no label</Info>
  error TS2741: Property 'label' is missing in type '{ children: string; }' but required in type 'InfoProps'.
<Box>{{ name: 'Suits' }}</Box>
  error TS2353: Object literal may only specify known properties, and 'name' does not exist in type 'ReactElement<...> | Iterable<ReactNode> | ReactPortal | Promise<...>'.
<WatchButton show="Suits" type="watch">Watch</WatchButton>
  error TS2322: Type '"watch"' is not assignable to type '"button" | "submit" | "reset" | undefined'.
<NoChildren title="Awake">text</NoChildren>
  error TS2322: Type '{ children: string; title: string; }' is not assignable to type 'IntrinsicAttributes & { title: string; }'.
  Property 'children' does not exist on type 'IntrinsicAttributes & { title: string; }'.
\end{lstlisting}

The plain object as children crashed \textbf{at run time} in JavaScript
(P3); TypeScript stops it \textbf{before} running. And a component whose
type has no \passthrough{\lstinline!children!} cannot be given content
at all.

\hypertarget{javascript-vs-typescript}{%
\subsection{JavaScript vs TypeScript}\label{javascript-vs-typescript}}

\begin{longtable}[]{@{}
  >{\raggedright\arraybackslash}p{(\columnwidth - 6\tabcolsep) * \real{0.2500}}
  >{\raggedright\arraybackslash}p{(\columnwidth - 6\tabcolsep) * \real{0.2500}}
  >{\raggedright\arraybackslash}p{(\columnwidth - 6\tabcolsep) * \real{0.2500}}
  >{\raggedright\arraybackslash}p{(\columnwidth - 6\tabcolsep) * \real{0.2500}}@{}}
\toprule\noalign{}
\begin{minipage}[b]{\linewidth}\raggedright
Form
\end{minipage} & \begin{minipage}[b]{\linewidth}\raggedright
Same in TypeScript
\end{minipage} & \begin{minipage}[b]{\linewidth}\raggedright
What I must add
\end{minipage} & \begin{minipage}[b]{\linewidth}\raggedright
Error when I get it wrong
\end{minipage} \\
\midrule\noalign{}
\endhead
\bottomrule\noalign{}
\endlastfoot
1 \passthrough{\lstinline!const [x, y] = f()!} & the pattern & a tuple
return type \passthrough{\lstinline!: [string, string]!} & without it
the items are \passthrough{\lstinline!string \\| undefined!}
(noUncheckedIndexedAccess) \\
2 \passthrough{\lstinline![, second, ...others]!} & skip and rest &
nothing; \passthrough{\lstinline!others!} is
\passthrough{\lstinline!string[]!} & using
\passthrough{\lstinline!second!} directly:
\passthrough{\lstinline!TS18048: 'second' is possibly 'undefined'!} \\
3 \passthrough{\lstinline![a, b] = [b, a]!} & the swap and the
\passthrough{\lstinline!;!} & \passthrough{\lstinline!let!} with a type
& none, the types must match \\
4 \passthrough{\lstinline![a = 1] = []!} & the default & nothing & none;
the default removes \passthrough{\lstinline!undefined!} from the type \\
5 \passthrough{\lstinline!\{ a, b \} = t!} & the pattern & a type for
the object (\passthrough{\lstinline!type ShowInfo = ...!}) & a field
that is not in the type:
\passthrough{\lstinline!TS2339: Property '...' does not exist on type!}
(T5) \\
6 \passthrough{\lstinline!\{ a: alpha \}!} & the rename & nothing;
trying to write a type here is the T1 mistake &
\passthrough{\lstinline!TS7031: Binding element 'string' implicitly has an 'any' type!} \\
7 \passthrough{\lstinline!\{ a = 1 \} = t!} & the default &
\passthrough{\lstinline!?!} on the field
(\passthrough{\lstinline!network?: string!}) & without
\passthrough{\lstinline!?!}, the default is useless and the object must
have the field \\
8 \passthrough{\lstinline!\{ a, ...others \}!} & rest and spread &
nothing; \passthrough{\lstinline!others!} gets the type without
\passthrough{\lstinline!a!} & none \\
9 nested & the pattern & a nested type
(\passthrough{\lstinline!owner: \{ firstName: string \}!}) & a wrong
path: \passthrough{\lstinline!TS2339: Property ... does not exist!} \\
10 \passthrough{\lstinline!\{ [key]: value \}!} & the computed key & a
type like \passthrough{\lstinline!Record<string, number>!} & the value
is \passthrough{\lstinline!number \\| undefined!}, a default makes it
\passthrough{\lstinline!number!} \\
11 parameters & the pattern and \passthrough{\lstinline!= \{\}!} & the
type \textbf{after the whole pattern}:
\passthrough{\lstinline!(\{ title \}: ShowOptions = \{\})!} & type
inside the pattern = rename, \passthrough{\lstinline!TS7031!} (T1) \\
12 loops & \passthrough{\lstinline!for (const [k, v] of ...)!} &
\passthrough{\lstinline!new Map<string, string>()!} & none, types come
from \passthrough{\lstinline!Object.entries!} / the Map \\
props & \passthrough{\lstinline!(\{ title, children \})!} &
\passthrough{\lstinline!children?: ReactNode!} or
\passthrough{\lstinline!PropsWithChildren!} &
\passthrough{\lstinline!Property 'children' does not exist on type 'IntrinsicAttributes \& ...'!} \\
\end{longtable}

\hypertarget{my-own-log}{%
\subsection{My own log}\label{my-own-log}}

Real things that went wrong or surprised me while I did this, in the
order they happened.

\begin{enumerate}
\def\labelenumi{\arabic{enumi}.}
\item
  \textbf{My Form 2 guesses were all wrong.} Before running it I wrote:
  \passthrough{\lstinline!secondShow!} → Suits,
  \passthrough{\lstinline!otherShows!} → Awake, and ``too short'' for
  the 4th item. The output was
  \passthrough{\lstinline!Awake [ 'The Mentalist' ]!} and
  \passthrough{\lstinline!undefined []!}. I did not know that the
  leading comma skips the first item, and I expected rest to give one
  value. Now I know rest always makes an array, even
  \passthrough{\lstinline![]!} when nothing is left, and a missing
  position is \passthrough{\lstinline!undefined!}, not an error.
\item
  \textbf{Form 3 without the semicolon.} I removed the
  \passthrough{\lstinline!;!} in front of
  \passthrough{\lstinline![favorite, secondFavorite] = [secondFavorite, favorite]!}
  and got
  \passthrough{\lstinline!ReferenceError: Cannot access 'secondFavorite' before initialization!}.
  JavaScript joined the line with the line above, so it read
  \passthrough{\lstinline!let secondFavorite = 'Awake'[favorite, secondFavorite] = ...!},
  which uses \passthrough{\lstinline!secondFavorite!} while creating it.
  A line that starts with \passthrough{\lstinline![!} or
  \passthrough{\lstinline!(!} needs a \passthrough{\lstinline!;!} in
  front when I do not write semicolons.
\item
  \textbf{Form 4: I guessed it would print ``empty, undefined, null''.}
  It printed \passthrough{\lstinline!Suits Awake!},
  \passthrough{\lstinline!The Mentalist Awake!} and
  \passthrough{\lstinline!null The Mentalist!}. The default replaced the
  missing item and \passthrough{\lstinline!undefined!}, but not
  \passthrough{\lstinline!null!}. The same thing came back in Form 7,
  Pitfall 2 and P6 (\passthrough{\lstinline!The Mentalist:!} with an
  empty space on the screen), so it is the rule I will remember most.
\item
  \textbf{Form 11 without \passthrough{\lstinline!= \{\}!}.} Calling
  \passthrough{\lstinline!describeShowNoDefault()!} gave
  \passthrough{\lstinline!Cannot destructure property 'title' of 'undefined' as it is undefined.!}
  I understood that a default for a field
  (\passthrough{\lstinline!seasons = 1!}) does not help when the whole
  argument is missing; the parameter itself needs
  \passthrough{\lstinline!= \{\}!}.
\item
  \textbf{T1: I tried to write types inside the pattern.}
  \passthrough{\lstinline!function showLabel(\{ name: string, seasons: number \})!}
  gave three errors:
  \passthrough{\lstinline!TS7031: Binding element 'string' implicitly has an 'any' type!}
  (twice) and
  \passthrough{\lstinline!TS2304: Cannot find name 'seasons'!}. The
  surprise was that \passthrough{\lstinline!name!} gave \textbf{no}
  error, even though it was renamed too. The reason: TypeScript knows
  the browser global \passthrough{\lstinline!name!}
  (\passthrough{\lstinline!window.name!}), so
  \passthrough{\lstinline!$\{name\}!} silently used that global instead.
  So this mistake can be hidden, not only loud.
\item
  \textbf{P2 and P3 surprises.} When I gave both
  \passthrough{\lstinline!children="from attribute"!} and content
  between the tags, the content between the tags won. In P3,
  \passthrough{\lstinline!false!} and \passthrough{\lstinline!null!}
  showed nothing but \passthrough{\lstinline!\{0\}!} printed
  \passthrough{\lstinline!0!}. I also saw that a function as children
  prints nothing and React warns \emph{Functions are not valid as a
  React child}, but it works when the component calls it.
\item
  \textbf{\passthrough{\lstinline!key!} and changing props.} Reading
  \passthrough{\lstinline!props.key!} gave
  \passthrough{\lstinline!undefined!} and React warned:
  \emph{\passthrough{\lstinline!key!} is not a prop}. Writing
  \passthrough{\lstinline!props.show = 'changed'!} threw
  \passthrough{\lstinline!TypeError - Cannot assign to read only property 'show'!}.
  I did not know React freezes props.
\item
  \textbf{The React warnings are not in
  \passthrough{\lstinline!output.txt!}.} PowerShell's
  \passthrough{\lstinline!>!} writes UTF-16, so I used
  \passthrough{\lstinline!npm run all | Out-File -Encoding utf8 output.txt!}.
  After that, the two React warnings still appeared in the terminal but
  not in the file, because React writes warnings to stderr and the pipe
  only takes stdout.
\item
  \textbf{Tooling problem.} To try the four P7 errors I pressed
  \passthrough{\lstinline!Ctrl + /!} to uncomment lines, but on my
  Turkish keyboard this zoomed VS Code instead and my selected lines
  were replaced. I used \emph{Edit → Toggle Line Comment} instead.
\end{enumerate}

\hypertarget{bonus-2-depth}{%
\subsection{Bonus 2: Depth}\label{bonus-2-depth}}

Bonus 2: Depth

All experiments are in \passthrough{\lstinline!js/depth.mjs!} (run with
\passthrough{\lstinline!node js/depth.mjs!}). Every output below is from
my run.

\hypertarget{d1-what-does-an-array-pattern-really-need}{%
\subsubsection{D1: what does an array pattern really
need?}\label{d1-what-does-an-array-pattern-really-need}}

My question: does \passthrough{\lstinline!const [a, b] = x!} need
\passthrough{\lstinline!x!} to be an array? Experiments:

\begin{lstlisting}
const [firstLetter, secondLetter, ...restLetters] = 'Suits'
const [firstGenre, ...otherGenres] = new Set(['drama', 'crime', 'drama', 'sci-fi'])
const [[firstShowInMap, firstLead]] = new Map([['Suits', 'Harvey'], ['The Mentalist', 'Patrick Jane']])
const [onlyFirst] = spyShows          // a generator that logs every next()
const [ep1, ep2, ep3] = episodeNumbers() // an infinite generator
const [x] = { name: 'Suits' }          // a plain object
const { 0: arrayFirst, length: showCount } = ['Suits', 'Awake', 'The Mentalist']
\end{lstlisting}

\begin{lstlisting}
D1 string: S u [ 'i', 't', 's' ]
D1 set (duplicates already gone): drama [ 'crime', 'sci-fi' ]
D1 map (first entry is itself an array): Suits Harvey
D1 spy: next() gave Suits
D1 spy result: Suits
D1 infinite generator: 1 2 3
D1 object with array pattern: TypeError - {(intermediate value)} is not iterable
D1 object pattern on an array: Suits 3
\end{lstlisting}

\textbf{What I learned:} an array pattern does not need an array, it
needs an \textbf{iterable}. It calls
\passthrough{\lstinline!value[Symbol.iterator]()!} and then
\passthrough{\lstinline!next()!} once per position. That is why a string
gives letters, a \passthrough{\lstinline!Set!} gives values without
duplicates, and a \passthrough{\lstinline!Map!} gives
\passthrough{\lstinline![key, value]!} pairs (so I needed a nested
pattern \passthrough{\lstinline![[show, lead]]!}). The spy shows
\passthrough{\lstinline!next()!} was called \textbf{only once} for
\passthrough{\lstinline![onlyFirst]!}: the pattern takes only as many
items as it needs, which is also why an \textbf{infinite} generator
works (with a rest element \passthrough{\lstinline!...rest!} it would
never finish). A plain object has no
\passthrough{\lstinline!Symbol.iterator!}, so it fails with \emph{is not
iterable}. The other direction works: an object pattern on an array
reads keys \passthrough{\lstinline!0!} and
\passthrough{\lstinline!length!}, because an array is an object.

\hypertarget{d2-order-of-evaluation-and-defaults}{%
\subsubsection{D2: order of evaluation and
defaults}\label{d2-order-of-evaluation-and-defaults}}

\begin{lstlisting}
function loudDefault(label) { console.log('D2 default computed for', label); return `default ${label}` }
const [given = loudDefault('given'), missing = loudDefault('missing')] = ['Suits']
const { title = 'Awake', subtitle = `${title} (2012)` } = {}
eval("const { first = second, second = 'x' } = {}")
\end{lstlisting}

\begin{lstlisting}
D2 default computed for missing
D2 values: Suits | default missing
D2 earlier variable in a default: Awake (2012)
D2 later variable in a default: ReferenceError - Cannot access 'second' before initialization
\end{lstlisting}

\textbf{What I learned:} a default is \textbf{lazy}: the function ran
only for \passthrough{\lstinline!missing!}, never for
\passthrough{\lstinline!given!}, because a value was there. So a default
with a side effect (a log, an API call) happens only sometimes. Patterns
run \textbf{left to right}: a default can use a variable that is earlier
in the same pattern (\passthrough{\lstinline!subtitle!} uses
\passthrough{\lstinline!title!}), but not a later one. That is the same
``before initialization'' error I got in Form 3 (log entry 2).

\hypertarget{d3-destructuring-in-other-places}{%
\subsubsection{D3: destructuring in other
places}\label{d3-destructuring-in-other-places}}

\begin{lstlisting}
try { JSON.parse('{ bad json') } catch ({ name, message }) { ... }
for (const key in seasonsByShow) ...
for (const [key, value] of Object.entries(seasonsByShow)) ...
;[ratings.suits, ratings.awake] = [9, 8]
;({ first: swapped.second, second: swapped.first } = { first: 'Suits', second: 'Awake' })
\end{lstlisting}

\begin{lstlisting}
D3 catch pattern: SyntaxError - Expected property name or '}' in JSON at position 2 (line 1 column 3)
D3 for...in key: Suits
D3 for...in key: Awake
D3 for...of pair: Suits 9
D3 for...of pair: Awake 1
D3 targets are properties: { suits: 9, awake: 8 }
D3 object pattern into properties: { second: 'Suits', first: 'Awake' }
\end{lstlisting}

\textbf{What I learned:} the \passthrough{\lstinline!catch!} parameter
is a normal parameter, so Form 5 works there.
\passthrough{\lstinline!for...in!} only gives \textbf{keys} as strings,
so there is nothing to destructure; \passthrough{\lstinline!for...of!}
with \passthrough{\lstinline!Object.entries!} gives pairs (Form 12). The
targets of an assignment pattern do not have to be variables: any place
I could write on the left of \passthrough{\lstinline!=!} works, also
object properties. In the last line the rename syntax wrote
\passthrough{\lstinline!first!} into
\passthrough{\lstinline!swapped.second!}, so the pattern also swapped
the names.

\hypertarget{d5-when-destructuring-makes-code-worse}{%
\subsubsection{D5: when destructuring makes code
worse}\label{d5-when-destructuring-makes-code-worse}}

\begin{lstlisting}
const { data: { show: { details: { cast: [{ actor: { name: leadActor } }] } } } } = apiResponse
// refactored:
const leadActorRefactored = apiResponse.data?.show?.details?.cast[0]?.actor?.name ?? 'unknown'
\end{lstlisting}

\begin{lstlisting}
D5 deep pattern: Simon Baker
D5 refactored: Simon Baker
D5 deep pattern on missing data: Cannot read properties of undefined (reading 'details')
\end{lstlisting}

\textbf{What I learned:} both give the same value, but the deep pattern
is hard to read (I had to count the braces), and it \textbf{crashes} as
soon as one level is missing. The refactored line reads left to right
like the path, and \passthrough{\lstinline!?.!} with
\passthrough{\lstinline!??!} gives a fallback instead of an error. My
rule now: destructure one or two levels, and use optional chaining for
long paths through data that may be incomplete.

\hypertarget{bonus-3-real-code}{%
\subsection{Bonus 3: Real code}\label{bonus-3-real-code}}

Bonus 3: Real code

Three sources, as the handout asks:

\begin{itemize}
\tightlist
\item
  \textbf{The Task Manager from Assignment 3} (my own copy of the
  starter project): cases A to C.
\item
  \textbf{Installed libraries in
  \passthrough{\lstinline!node\_modules!}}:
  \passthrough{\lstinline!tsx!} and
  \passthrough{\lstinline!@types/react!} (DefinitelyTyped): cases 2 to
  6.
\item
  \textbf{An open-source project of my choice: the TypeScript compiler}
  (Microsoft, Apache-2.0, installed here as
  \passthrough{\lstinline!typescript!} 5.8.3): case 1.
\end{itemize}

File and line numbers are as in my files.

\hypertarget{from-the-task-manager-assignment-3}{%
\subsubsection{From the Task Manager (Assignment
3)}\label{from-the-task-manager-assignment-3}}

\textbf{Case A --- Form 1 with a React hook.}
\passthrough{\lstinline!assignment3/src/App.tsx!}, line 72:

\begin{lstlisting}
const [state, dispatch] = useReducer(taskReducer, initialState)
\end{lstlisting}

Without destructuring:
\passthrough{\lstinline!const pair = useReducer(taskReducer, initialState); const state = pair[0]; const dispatch = pair[1];!}
\passthrough{\lstinline!useReducer!} returns a tuple
\passthrough{\lstinline![state, dispatch]!}. This is the line the
handout mentions; I wrote it in Assignment 3 without thinking of it as
destructuring.

\textbf{Case B --- React props with no children.}
\passthrough{\lstinline!assignment3/src/TaskRow.tsx!}, line 11:

\begin{lstlisting}
export default function TaskRow({ task, onNext, onDelete }: TaskRowProps) {
\end{lstlisting}

used in \passthrough{\lstinline!App.tsx!}, lines 60 to 65:

\begin{lstlisting}
<TaskRow key={t.id} task={t} onNext={() => dispatch({ type: 'advance', id: t.id })}
         onDelete={() => dispatch({ type: 'remove', id: t.id })} />
\end{lstlisting}

Without destructuring:
\passthrough{\lstinline!function TaskRow(props: TaskRowProps) \{ ... props.task.title ... props.onNext ... \}!}
\textbf{What is a prop and what is children here:}
\passthrough{\lstinline!task!}, \passthrough{\lstinline!onNext!} and
\passthrough{\lstinline!onDelete!} are props (attributes in the JSX).
There are \textbf{no children}:
\passthrough{\lstinline!<TaskRow ... />!} is self-closing and
\passthrough{\lstinline!TaskRowProps!} has no
\passthrough{\lstinline!children!} field, so giving it content would be
the TS2322 error from P7. \passthrough{\lstinline!key=\{t.id\}!} is
written like a prop, but as I saw in P4 React keeps it and
\passthrough{\lstinline!TaskRow!} never receives it. This is Form 5 in a
parameter (Form 11) with the type after the whole pattern (T1).

\textbf{Case C --- Form 11 with an inline type.}
\passthrough{\lstinline!assignment3/src/App.tsx!}, line 8:

\begin{lstlisting}
function AddForm({ dispatch }: { dispatch: Dispatch<Action> }) {
\end{lstlisting}

Without destructuring:
\passthrough{\lstinline!function AddForm(props: \{ dispatch: Dispatch<Action> \}) \{ const dispatch = props.dispatch; ... \}!}
The component only needs one prop, so the type is written inline after
the pattern. If I wrote
\passthrough{\lstinline!\{ dispatch: Dispatch<Action> \}!} \emph{inside}
the pattern, it would be the T1 rename mistake.

\hypertarget{from-installed-libraries-and-the-typescript-project}{%
\subsubsection{From installed libraries and the TypeScript
project}\label{from-installed-libraries-and-the-typescript-project}}

\textbf{Case 1 --- Form 5, a real Node.js idiom.}
\passthrough{\lstinline!node\_modules/typescript/lib/tsc.js!}, line 3:

\begin{lstlisting}
const { enableCompileCache } = require("node:module");
\end{lstlisting}

Without destructuring:
\passthrough{\lstinline!const enableCompileCache = require("node:module").enableCompileCache;!}
TypeScript takes one function out of a module. The next line checks
\passthrough{\lstinline!if (enableCompileCache)!}, because on older Node
versions the field is simply missing and the pattern gives
\passthrough{\lstinline!undefined!} instead of an error (Form 5: a
missing field is \passthrough{\lstinline!undefined!}).

\textbf{Case 2 --- Form 1, the most common destructuring in React.}
\passthrough{\lstinline!node\_modules/@types/react/index.d.ts!}, line
1640 (in a doc example):

\begin{lstlisting}
const [count, setCount] = useState(0);
\end{lstlisting}

Without destructuring:
\passthrough{\lstinline!const state = useState(0); const count = state[0]; const setCount = state[1];!}
\passthrough{\lstinline!useState!} returns a tuple
\passthrough{\lstinline![value, setter]!}, so the names are chosen by
position. This is exactly why React returns an array and not an object:
with an array I can name the two values anything I like without a rename
(Form 6).

\textbf{Case 3 --- React props and children (Form 5 + Form 11).}
\passthrough{\lstinline!node\_modules/@types/react/index.d.ts!}, line
419:

\begin{lstlisting}
type Props = { children: ReactNode }
const Component = ({ children }: Props) => <div>{children}</div>
<Component>hello</Component>
\end{lstlisting}

Without destructuring:
\passthrough{\lstinline!const Component = (props: Props) => <div>\{props.children\}</div>!}
Here there are \textbf{no other props}: the only thing in the props
object is \passthrough{\lstinline!children!}, and its value is the text
\passthrough{\lstinline!hello!} written between the tags. The type is
written after the whole pattern (T1), and
\passthrough{\lstinline!children!} is required (no
\passthrough{\lstinline!?!}), so \passthrough{\lstinline!<Component />!}
without content would be a compile error.

\textbf{Case 4 --- Form 1 on a function result, and a case where it
makes the code worse.}
\passthrough{\lstinline!node\_modules/tsx/dist/cli.mjs!}, line 22:

\begin{lstlisting}
[W,p]=process.versions.node.split(".").slice(0,2).map(m=>Number.parseInt(m,10))
\end{lstlisting}

Without destructuring:
\passthrough{\lstinline!const parts = process.versions.node.split(".").map(Number); const W = parts[0]; const p = parts[1];!}
It reads the Node major and minor version (for me
\passthrough{\lstinline!24!} and \passthrough{\lstinline!11!}).
\textbf{Why I think this is worse:} this file is minified, so the names
\passthrough{\lstinline!W!} and \passthrough{\lstinline!p!} say nothing,
and the pattern hides that the two values mean ``major'' and ``minor''.
In source code \passthrough{\lstinline!const [major, minor] = ...!}
would be clear; with one-letter names the position is the only meaning,
and swapping them would not give any error (Pitfall 5).

\textbf{Case 5 --- Form 6 (rename) used by a minifier.}
\passthrough{\lstinline!node\_modules/tsx/dist/cli.mjs!}, line 2:

\begin{lstlisting}
{stringify:Be}=JSON, ... {hasOwnProperty:Pr}=Object.prototype
\end{lstlisting}

Without destructuring:
\passthrough{\lstinline!const Be = JSON.stringify; const Pr = Object.prototype.hasOwnProperty;!}
The build tool takes methods out of built-in objects once and renames
them to short names. Later the code calls
\passthrough{\lstinline!Pr.call(e,t)!}: because
\passthrough{\lstinline!hasOwnProperty!} was taken out of the object, it
lost its \passthrough{\lstinline!this!}, so it has to be called with
\passthrough{\lstinline!.call!}. This is a real trap of destructuring
methods: the function is copied, not bound to the object.

\textbf{Case 6 --- Form 12, nested loops over entries.}
\passthrough{\lstinline!node\_modules/tsx/dist/cli.mjs!}, line 3:

\begin{lstlisting}
for(let[u,n]of Object.entries(t)){for(let[s,r]of Object.entries(n)) ... }
\end{lstlisting}

Without destructuring:
\passthrough{\lstinline!for (const entry of Object.entries(t)) \{ const u = entry[0]; const n = entry[1]; ... \}!}
The terminal colour table is an object of objects, so the code loops
over the outer pairs and then over the inner pairs, destructuring each
\passthrough{\lstinline![key, value]!} pair, exactly like my Form 12
example.

\hypertarget{bonus-4-git-repository}{%
\subsection{Bonus 4: Git repository}\label{bonus-4-git-repository}}

Bonus 4: Git repository ---
https://github.com/ranakurts/sen0415-assignment4

The repository has a \passthrough{\lstinline!.gitignore!} that leaves
out \passthrough{\lstinline!node\_modules!}, and it runs after
\passthrough{\lstinline!git clone!},
\passthrough{\lstinline!npm install!} and
\passthrough{\lstinline!npm run all!}. The history shows the work in
steps (\passthrough{\lstinline!git log --oneline!}): starter, my name,
array forms, object forms, parameter and loop forms, pitfalls, props
P1--P3 and P4--P6, the TypeScript forms, T1--T5, typed props P1--P7,
\passthrough{\lstinline!output.txt!}, the README, and each bonus as its
own commit. Honest note: I did all commits on the same evening (20
November), so the history is in many steps but not over several days.

\hypertarget{bonus-1-latex}{%
\subsection{Bonus 1: LaTeX}\label{bonus-1-latex}}

Bonus 1: LaTeX

The guide is also in \passthrough{\lstinline!latex/guide.tex!} with the
compiled \passthrough{\lstinline!latex/guide.pdf!}. It has the same
content as this README (all twelve entries, the pitfalls, T1--T5, props
and children, the table, my log and the bonuses), a contents page with
links to every section, code in \passthrough{\lstinline!listings!} boxes
and real tables.

Commands I used:

\begin{lstlisting}
sed '1d' README.md > body.md     # drop the "# Destructuring Field Guide" line, the title page has it
pandoc body.md -s --listings --toc -H latex/header.tex -V geometry:margin=2.2cm -V fontsize=10pt \
  -V colorlinks=true -V title="Destructuring Field Guide" -V author="Rana Neslihan Kurt, 2300006496" \
  -V date="SEN0415 Advanced Web Programming, Assignment 4" -o latex/guide.tex
cd latex
pdflatex guide.tex
pdflatex guide.tex
\end{lstlisting}

(\passthrough{\lstinline!pdflatex!} twice so that the contents page and
the links are correct.) The build finished with no errors and no
overfull boxes.

\end{document}
